% TOP_PROBLEM: {"id": 30006829, "title": "Michel Boundary Rigidity Conjecture for Simple Manifolds", "category_id": 6, "category": {"id": 6, "name": "geometry", "display_name": "Geometry", "description": "Euclidean and non-Euclidean geometry, geometric structures.", "slug": "geometry", "order_index": 6, "created_at": "2026-07-31T15:26:25.671Z"}, "status": "open"}
% FIRST_POSTED: 2026-09-25
% ENTRY_KIND: statement_only
% !TeX program = lualatex
% Definition and sources only; this file is not an LLM solution attempt.
\documentclass[11pt]{article}
\usepackage[margin=25mm]{geometry}
\usepackage{fontspec}
\setmainfont{Latin Modern Roman}
\usepackage{amsmath,amssymb,mathtools}
\usepackage{unicode-math}
\setmathfont{Latin Modern Math}
\usepackage{xurl,hyperref}
\hypersetup{colorlinks=true,urlcolor=blue}
\setcounter{secnumdepth}{0}
\setlength{\parindent}{0pt}
\setlength{\parskip}{5pt}
\setlength{\emergencystretch}{3em}
\begin{document}
\section*{\texorpdfstring{Michel Boundary Rigidity Conjecture for Simple Manifolds}{Michel Boundary Rigidity Conjecture for Simple Manifolds}}\label{30006829}
\subsection{Definitions and mathematical statement}
A simple Riemannian metric on a compact manifold with boundary has strictly convex boundary, no conjugate points, and unique geodesic connections between points, in the cited convention. Its boundary distance function is
$d_g:\partial M\times\partial M\to{\mathbb{R}}_{\ge0}$, computed using curves in the whole manifold. For smooth simple metrics $g,g'$ on the same $n$-manifold, $n\ge3$, the conjecture asserts
\[
 d_g|_{\partial M\times\partial M}=d_{g'}|_{\partial M\times\partial M}
 \quad\Longrightarrow\quad g'=\phi^*g,
 \qquad\phi|_{\partial M}=\mathrm{id}.
\]
The coordinate ambiguity is unavoidable. Boundary-to-boundary distances are the given data, not all interior distances or a full Dirichlet-to-Neumann operator.
\par\smallskip\textit{Formulation and concept references:} \href{https://arxiv.org/pdf/1702.03638v3}{[S1]}.

\subsection{Short English statement}
Do all boundary-to-boundary distances determine a simple interior Riemannian metric, up to a diffeomorphism fixing the boundary?
\subsection{Sources}
\begin{itemize}
\item[S1]Local and global boundary rigidity and the geodesic X-ray transform in the normal gauge, arXiv1702.03638v3,12May2021; internal revision13May2021.
\url{https://arxiv.org/pdf/1702.03638v3}.
\textit{Formulation source.}
\end{itemize}
\end{document}
